\par\Needspace{6\baselineskip}\medskip
\hypertarget{native-141-article101}{}\pdfbookmark[2]{101.}{article-101}
\begin{center}101.\end{center}
If we introduce the quantity $\zeta$ and set $\BareRoot{(l-z)}=\dfrac my$ or $\BareRoot{(L+z)}=\dfrac MY$, and also set
\[
\mathrm{15)}\qquad\frac{mm}{\tfrac56+l+\zeta}=h
\]
\noindent or
\[
\mathrm{15^*)}\qquad\frac{MM}{L-\tfrac56-\zeta}=H,
\]
\noindent then equations 13 and $13^*$ take the form
\[
\begin{aligned}
\mathrm{16)}\qquad&\frac{(y-1)yy}{y+\tfrac19}=h\\[.6em]
\mathrm{16^*)}\qquad&\frac{(Y+1)YY}{Y-\tfrac19}=H,
\end{aligned}
\]
\noindent and thus become identical to those obtained for the ellipse (15 and $15^*$, Art. 91). Hence, insofar as $h$ or $H$ can be treated as known, $y$ or $Y$ can be deduced, and we then have
\[
\begin{aligned}
\mathrm{17)}\qquad&z=l-\frac{mm}{yy}\\[.6em]
\mathrm{17^*)}\qquad&z=\frac{MM}{YY}-L.
\end{aligned}
\]
\noindent It follows that all the operations prescribed above for the ellipse apply equally to the hyperbola, up to the step of using the approximate value of $h$ or $H$ to obtain $y$ or $Y$. The quantity $\dfrac{mm}{yy}-l$ or $L-\dfrac{MM}{YY}$, which had to be positive for the ellipse and $=0$ for the parabola, must then be negative for the hyperbola. This criterion therefore determines the kind of conic section. From the value of $z$ thus found, our table gives $\zeta$, which yields a corrected value of $h$ or $H$. The calculation must be repeated with this value until everything agrees exactly.
\par
Once the true value of $z$ has been found, $c$ could be deduced from it by the formula $c=1+2z+2\BareRoot{(z+zz)}$. With the applications that follow also in view, it is preferable to introduce the auxiliary angle $n$, to be determined by the equation $\operatorname{tang}2n=2\BareRoot{(z+zz)}$. This gives $c=\operatorname{tang}2n+\BareRoot{(1+\operatorname{tang}2n^2)}=\operatorname{tang}(45^0+n)$.
